\clearpage\thispagestyle{empty}
\pdfbookmark[0]{Quellen zu Beobachtungen der Ceres in Göttingen (1829)}{ceres1829-sources}
\begin{center}\textbf{Quellen zu Beobachtungen der Ceres in Göttingen (1829)}\end{center}
\noindent Diese KI-Übertragung ins heutige Deutsch folgt dem Beitrag auf S.~462 der \href{https://gdz.sub.uni-goettingen.de/id/PPN236007467}{\textit{Werke}, Band~VI} (Göttingen, 1874). Carl Friedrich Gauß’ Beobachtungen erschienen erstmals als \textit{Beobachtungen der Ceres in Göttingen} in den \textit{Astronomischen Nachrichten}, Band~VII, Nr.~161, \href{https://api.digitale-sammlungen.de/iiif/image/v2/bsb10049432_00211/full/1800,/0/default.jpg}{Sp.~329--330}, im Juni 1829.
\par\smallskip\noindent Die Werkausgabe nennt Gauß in der eingeklammerten Überschrift und fasst alle acht Beobachtungen in einer Tabelle zusammen. Das Feld für die Deklination am 16. März bleibt dort leer. Der Erstdruck ordnet die Beobachtungen in zwei Blöcken an, setzt einen Strich in das Feld und schließt mit Gauß’ Unterschrift. KI-Übertragung und KI-Satz.
\clearpage
